% GENERATED FILE. Edit canonical lessons, then run npm run generate:books.
% canonical-source-hash: fnv1a64:3d48399a393e97dd
% canonical-lessons: IT-C292-rapina, IT-C292-rapinare, IT-C292-rapinatore, IT-C292-cassaforte, IT-C292-bottino

\chapter{Rapina in banca}
\label{ch:it-rapina-in-banca}
\hlchaptermodality{292}

\begin{chapteropening}
\textbf{By the end of this chapter:} \emph{I can report a robbery: robbing, the robbers, the safe and the loot.}

Complete the last lesson of chapter 292: I can report a robbery: robbing, the robbers, the safe and the loot.
\end{chapteropening}

\section[la rapina]{la rapina — the robbery, the hold-up}

\label{lesson:IT-C292-rapina}



\begin{quote}
Before the new one: say the Italian for to rob someone, then the Italian for the stolen goods.
\end{quote}

\subsection*{You'll want to know: la rapina}
\textbf{la rapina} — \textquotedblleft{}the robbery, the hold-up\textquotedblright{}.

\begin{quote}
C'è stata una rapina in banca stamattina.
\end{quote}

\begin{quote}
\emph{There was a bank robbery this morning.}
\end{quote}

With force or threats; \emph{a mano armata} means armed.

\begin{cousinweb}
\emph{rapina} is inherited from Latin \emph{rapīna} \textquotedblleft{}robbery, plunder\textquotedblright{} (EtymDB), from \emph{rapere} \textquotedblleft{}to seize\textquotedblright{}.
\end{cousinweb}

\begin{grammarlens}[title={What you've built}]
One more word for telling, explaining and describing at B1.
\end{grammarlens}

\subsection*{Your turn}
\begin{itemize}
\raggedright
  \item \emph{Say it:} \emph{la rapina}
  \item \emph{Say it:} \emph{la rapina}, once more
  \item \emph{Say it:} a sentence of your own with \emph{la rapina}
  \item \emph{Recall:} say the Italian for to rob someone, then the Italian for the stolen goods, then say \emph{la rapina} again
\end{itemize}

\begin{tcolorbox}[breakable,colback=teal!4,colframe=teal!35!black,title={Before you move on}]
What does la rapina mean? (\textquotedblleft{}The robbery, the hold-up\textquotedblright{}.) Say it once more.
\end{tcolorbox}
\section[rapinare]{rapinare — to rob (with force)}

\label{lesson:IT-C292-rapinare}



\begin{quote}
Before the new one: say the Italian for the stolen goods, then the Italian for the robbery.
\end{quote}

\subsection*{You'll want to know: rapinare}
\textbf{rapinare} — \textquotedblleft{}to rob (with force)\textquotedblright{}.

\begin{quote}
Due uomini hanno rapinato il supermercato.
\end{quote}

\begin{quote}
\emph{Two men robbed the supermarket.}
\end{quote}

\emph{Rapinare una banca}; \emph{rapinare qualcuno}.

\begin{cousinweb}
\emph{rapinare} is built on \emph{rapina} \textquotedblleft{}robbery\textquotedblright{}, from Latin \emph{rapīna} and \emph{rapere} \textquotedblleft{}to seize\textquotedblright{} (EtymDB).
\end{cousinweb}

\begin{grammarlens}[title={What you've built}]
One more word for telling, explaining and describing at B1.
\end{grammarlens}

\subsection*{Your turn}
\begin{itemize}
\raggedright
  \item \emph{Say it:} \emph{rapinare}
  \item \emph{Say it:} \emph{rapinare}, once more
  \item \emph{Say it:} a sentence of your own with \emph{rapinare}
  \item \emph{Recall:} say the Italian for the stolen goods, then the Italian for the robbery, then say \emph{rapinare} again
\end{itemize}

\begin{tcolorbox}[breakable,colback=teal!4,colframe=teal!35!black,title={Before you move on}]
What does rapinare mean? (\textquotedblleft{}To rob (with force)\textquotedblright{}.) Say it once more.
\end{tcolorbox}
\section[il rapinatore]{il rapinatore — the robber}

\label{lesson:IT-C292-rapinatore}



\begin{quote}
Before the new one: say the Italian for the robbery, then the Italian for to rob.
\end{quote}

\subsection*{You'll want to know: il rapinatore}
\textbf{il rapinatore} — \textquotedblleft{}the robber\textquotedblright{}.

\begin{quote}
Il rapinatore è scappato in moto.
\end{quote}

\begin{quote}
\emph{The robber got away on a motorbike.}
\end{quote}

\emph{-tore} makes the person who does it: \emph{rapinare}, \emph{rapinatore}.

\begin{cousinweb}
\emph{rapinatore} is \emph{rapinare} with the agent ending \emph{-tore}, from Latin \emph{-tor}; the root is Latin \emph{rapere} \textquotedblleft{}to seize\textquotedblright{} (EtymDB).
\end{cousinweb}

\begin{grammarlens}[title={What you've built}]
One more word for telling, explaining and describing at B1.
\end{grammarlens}

\subsection*{Your turn}
\begin{itemize}
\raggedright
  \item \emph{Say it:} \emph{il rapinatore}
  \item \emph{Say it:} \emph{il rapinatore}, once more
  \item \emph{Say it:} a sentence of your own with \emph{il rapinatore}
  \item \emph{Recall:} say the Italian for the robbery, then the Italian for to rob, then say \emph{il rapinatore} again
\end{itemize}

\begin{tcolorbox}[breakable,colback=teal!4,colframe=teal!35!black,title={Before you move on}]
What does il rapinatore mean? (\textquotedblleft{}The robber\textquotedblright{}.) Say it once more.
\end{tcolorbox}
\section[la cassaforte]{la cassaforte — the safe}

\label{lesson:IT-C292-cassaforte}



\begin{quote}
Before the new one: say the Italian for to rob, then the Italian for the robber.
\end{quote}

\subsection*{You'll want to know: la cassaforte}
\textbf{la cassaforte} — \textquotedblleft{}the safe\textquotedblright{}.

\begin{quote}
I soldi erano nella cassaforte.
\end{quote}

\begin{quote}
\emph{The money was in the safe.}
\end{quote}

Plural \emph{le casseforti}.

\begin{cousinweb}
\emph{cassaforte} joins \emph{cassa} \textquotedblleft{}box\textquotedblright{}, from Latin \emph{capsa} \textquotedblleft{}a case\textquotedblright{}, and \emph{forte} \textquotedblleft{}strong\textquotedblright{}, from Latin \emph{fortis} (Lewis \& Short).
\end{cousinweb}

\begin{grammarlens}[title={What you've built}]
One more word for telling, explaining and describing at B1.
\end{grammarlens}

\subsection*{Your turn}
\begin{itemize}
\raggedright
  \item \emph{Say it:} \emph{la cassaforte}
  \item \emph{Say it:} \emph{la cassaforte}, once more
  \item \emph{Say it:} a sentence of your own with \emph{la cassaforte}
  \item \emph{Recall:} say the Italian for to rob, then the Italian for the robber, then say \emph{la cassaforte} again
\end{itemize}

\begin{tcolorbox}[breakable,colback=teal!4,colframe=teal!35!black,title={Before you move on}]
What does la cassaforte mean? (\textquotedblleft{}The safe\textquotedblright{}.) Say it once more.
\end{tcolorbox}
\section[il bottino]{il bottino — the loot, the haul}

\label{lesson:IT-C292-bottino}



\begin{quote}
Before the new one: say the Italian for the robber, then the Italian for the safe.
\end{quote}

\subsection*{You'll want to know: il bottino}
\textbf{il bottino} — \textquotedblleft{}the loot, the haul\textquotedblright{}.

\begin{quote}
Il bottino è di diecimila euro.
\end{quote}

\begin{quote}
\emph{The haul is ten thousand euros.}
\end{quote}

Also the spoils of war.

\begin{cousinweb}
\emph{bottino} was borrowed from French \emph{butin} \textquotedblleft{}booty\textquotedblright{} (EtymDB), which goes back to a Germanic word for booty, related to German \emph{Beute} (EtymDB).
\end{cousinweb}

\begin{grammarlens}[title={What you've built}]
One more word for telling, explaining and describing at B1.
\end{grammarlens}

\subsection*{Your turn}
\begin{itemize}
\raggedright
  \item \emph{Say it:} \emph{il bottino}
  \item \emph{Say it:} \emph{il bottino}, once more
  \item \emph{Say it:} a sentence of your own with \emph{il bottino}
  \item \emph{Recall:} say the Italian for the robber, then the Italian for the safe, then say \emph{il bottino} again
\end{itemize}

\begin{tcolorbox}[breakable,colback=teal!4,colframe=teal!35!black,title={Before you move on}]
What does il bottino mean? (\textquotedblleft{}The loot, the haul\textquotedblright{}.) Say it once more.
\end{tcolorbox}
